% This is text associated with the open textbook "Open Electromagnetics"
% See immediately below for author, date
% This work is licensed under the Creative Commons Attribution-ShareAlike 4.0 International License. (CC BY-SA 4.0) 
% To view a copy of the license, visit http://creativecommons.org/licenses/by-sa/4.0/

%ORB TITLE Why 50 Ohms?
%ORB AUTHOR 0        # S.W. Ellingson, Virginia Tech (ellingson@vt.edu)
%ORB DATE 20191202   # yyyymmdd
%ORB ABSTRACT  

\label{m0191_Why_50_Ohms} % <-- Must match name of this file and module directory name.  Don't mess with this!

The quantity $50~\Omega$ appears in a broad range of applications across the field of electrical engineering.
In particular, it is a very popular value for the characteristic impedance of transmission line, and is commonly specified as the port impedance for signal sources, amplifiers, filters, antennas, and other RF components. 
So, what's special about $50~\Omega$?
The short answer is ``nothing.''
In fact, other standard impedances are in common use -- prominent among these is $75~\Omega$.
It is shown in this section that a broad range of impedances -- on the order of 10s of ohms -- emerge as useful values based on technical considerations such as minimizing attenuation, maximizing power handling, and compatibility with common types of antennas.
Characteristic impedances up to $300~\Omega$ and beyond are useful in particular applications.
However, it is not practical or efficient to manufacture and sell products for every possible impedance in this range.
Instead, engineers have settled on $50~\Omega$ as a round number that lies near the middle of this range, and have chosen a few other values to accommodate the smaller number of applications where there may be specific compelling considerations.
 
So, the question becomes ``what makes characteristic impedances in the range of 10s of ohms particularly useful?''
One consideration is attenuation 
\index{attenuation}
in coaxial cable.  
\index{transmission line!coaxial}
Coaxial cable is by far the most popular type of transmission line for connecting devices on separate printed circuit boards or in separate enclosures.
The attenuation of coaxial cable is addressed in Section~\ref{m0189_Attenuation_Coaxial_Cable}.
In that section, it is shown that attenuation is minimized for characteristic impedances in the range $\left(60~\Omega\right)/\sqrt{\epsilon_r}$ to 
$\left(77~\Omega\right)/\sqrt{\epsilon_r}$, where $\epsilon_r$ is the relative permittivity of the spacer material.
So, we find that $Z_0$ in the range $60~\Omega$ to $77~\Omega$ is optimum for air-filled cable, but more like $40~\Omega$ to $50~\Omega$ for cables using a plastic spacer material having typical $\epsilon_r\approx 2.25$.
Thus, $50~\Omega$ is clearly a reasonable choice if a single standard value is to be established for all such cable.
   
Coaxial cables are often required to carry high power signals.  In such applications, power handling
\index{transmission line!power handling}
capability is also important, and is addressed in Section~\ref{m0190_Power_Handling_Coaxial_Cable}.
In that section, we find the power handling capability of coaxial cable is optimized when the ratio of radii of the outer to inner conductors $b/a$ is about 1.65.  
For the air-filled cables typically used in high-power applications, this corresponds to a characteristic impedance of about $30~\Omega$.
This is significantly less than the $60~\Omega$ to $77~\Omega$ that minimizes attenuation in air-filled cables.
So, $50~\Omega$ can be viewed as a compromise between minimizing attenuation and maximizing power handling in air-filled coaxial cables. 

Although the preceding arguments justify $50~\Omega$ as a standard value, one can also see how one might make a case for $75~\Omega$ as a secondary standard value, especially for applications where attenuation is the primary consideration.  

Values of $50~\Omega$ and $75~\Omega$ also offer some convenience when connecting RF devices to antennas.  
For example, $75~\Omega$ is very close to the impedance of the commonly-encountered half-wave dipole antenna 
\index{antenna!half-wave dipole}
(about $73+j42~\Omega$), which may make impedance matching to that antenna easier. 
Another commonly-encountered antenna is the quarter-wave monopole, 
\index{antenna!quarter-wave monopole}
which exhibits an impedance of about $36+j21~\Omega$, which is close to $50~\Omega$.
In fact, we see that if we desire a single characteristic impedance that is equally convenient for applications involving either type of antenna, then $50~\Omega$ is a reasonable choice.

A third commonly-encountered antenna is the \emph{folded} half-wave dipole.  
\index{antenna!folded half-wave dipole}
This type of antenna is similar to a half-wave dipole but has better bandwidth, and is commonly used in FM and TV systems and land mobile radio (LMR) base stations.
A folded half-wave dipole has an impedance of about $300~\Omega$ and is balanced (not single-ended); 
\index{transmission line!single-ended}
thus, there is a market for balanced transmission line having $Z_0=300~\Omega$.  
However, it is very easy and inexpensive to implement a balun 
\index{balun}
(a device which converts the dipole output from balanced to unbalanced) while simultaneously stepping down impedance by a factor of 4; i.e., to $75~\Omega$.
Thus, we have an additional application for $75~\Omega$ coaxial line.

Finally, note that it is quite simple to implement microstrip 
\index{microstrip}
transmission line having characteristic impedance in the range $30~\Omega$ to $75~\Omega$.  For example, $50~\Omega$ on commonly-used 1.575~mm FR4 
\index{FR4}
requires a width-to-height ratio of about 2, so the trace is about 3~mm wide.  This is a very manageable size and easily implemented in printed circuit board designs.

%------------------------------------

{\bf Additional Reading:}
\begin{itemize}
\item ``\href{https://en.wikipedia.org/wiki/Dipole_antenna}{Dipole antenna}'' on Wikipedia.
\item ``\href{https://en.wikipedia.org/wiki/Monopole_antenna}{Monopole antenna}'' on Wikipedia.
\item ``\href{https://en.wikipedia.org/wiki/Balun}{Balun}'' on Wikipedia.
\end{itemize}

